% CONFIGURAÇÕES DE APARÃÅ NCIA DO PDF FINAL--------------------------------------
\makeatletter
\hypersetup{%
    portuguese,
    colorlinks=true,   % true: "links" coloridos; false: "links" em caixas de texto
    linkcolor=black,   % Links internos em preto (padrão do modelo UTFPR, sem vínculos coloridos)
    citecolor=black,   % Links de citações/referências em preto
    filecolor=black,   % Links de arquivos em preto
    urlcolor=black,    % URLs em preto
    breaklinks=true,
    pdftitle={\@title},
    pdfauthor={\@author},
    pdfkeywords={RAIS, CAGED, Mercado de TI, Engenharia de Dados, Arquitetura Medalhão, DuckDB, MinIO, Streamlit, Prophet, Forecasting}
}
\makeatother

% ALTERA O ASPECTO DA COR AZUL--------------------------------------------------
\definecolor{blue}{RGB}{41,5,195}

% REDEFINIÇÃÆ’O DE LABELS---------------------------------------------------------
\def\equationautorefname~#1\null{Equa\c c\~ao~(#1)\null}

% CRIA ÍNDICE REMISSIVO---------------------------------------------------------
\makeindex

% "Citado na página X" nas referências (gerado pelo abnt-refinfo do abntex2cite)
% A NBR 6023 não prevê essa informação, e os TCCs da UTFPR não a trazem: suprimida.
\renewcommand*{\backref}[1]{}
\renewcommand*{\backrefalt}[4]{}

% URL sem < > nas referências, como pede a NBR 6023:2018 (o abntex2cite ainda os põe).
\AtBeginDocument{\def\UrlLeft{}\def\UrlRight{}}

% CABEÇALHO DAS PÁGINAS---------------------------------------------------------
% Só o número da página, no canto superior direito, sem título do capítulo e
% sem filete (NBR 14724), como nos TCCs da UTFPR. Anexado a \textual porque é
% lá que a classe liga o estilo abntheadings.
\makepagestyle{numeropagina}
\makeoddhead{numeropagina}{}{}{\ABNTEXfontereduzida\thepage}
\makeevenhead{numeropagina}{}{}{\ABNTEXfontereduzida\thepage}
\makeheadrule{numeropagina}{\textwidth}{0pt}
\makeatletter
\g@addto@macro\textual{\pagestyle{numeropagina}}
\makeatother

% HIFENIZAÇÃÆ’O DE PALAVRAS QUE NÃÆ’O ESTÃÆ’O NO DICIONÁRIO---------------------------
\hyphenation{%
    qua-dros-cha-ve
    Kat-sa-gge-los
}

